\section{Conclusion}
\label{sec:conclusion}

Where to spend a per-epoch budget $B$ to recover recent top-$10$ neighbors depends on vector placement and topic drift. Under centroid placement, setup summaries already guide queries to shards receiving new vectors, so full refresh every epoch changes recent recall@10 over never refreshing by only \RNumHOneGapSteady{} on steady and \RNumHOneGapNewTopic{} on new-topic. At $Q = 200, B = 4$, causal online extra search and vector push beat every refresh rule on both streams. Data-change refresh overtakes extra search only on new-topic at $Q = 800$, by \RNumQSweepRefMinusESQEightHundredNewTopicAbsMean{} uncharged and \RNumAddThreeXFourESNewTopicAbsMean{} with its score bytes charged (fragile, not ahead on overall recall@10), while trailing vector push. Under arrival-time placement, \RNumCentersPerSummary{}-center summaries fail to summarize random steady arrivals and change overall recall@10 by \RNumGateMTwoOverallSteady{} and \RNumGateMThreeOverallSteady{} when refreshed, leaving vector push ahead at every $B$. On new-topic arrivals, refreshing active summaries raises recent recall@10 by \RNumGateMTwoNewTopic{} in single-segment placement and \RNumGateMThreeNewTopic{} in eight-segment placement, beating vector push across $B \in \{2, 4, 8\}$ in single-segment placement. In eight-segment placement on new-topic, refresh has no clear winner against vector push at $B = 4$ and loses at $B = 8$. At $B = 2$ it edges vector push only with score bytes uncharged; charging them costs a whole summary under the zero-slack floor and leaves no lead (\RNumAddThreeXTwoPVBTwoNewTopic{}). There, at $B = 4$, recovered-neighbor loses to data-change by \RNumHThreeBRNMinusDCNewTopic{} with score bytes uncharged, mainly because the probe picks worse shards. With data-change's \RNumDCScoreEpochBytes{} score bytes/epoch charged so both get \RNumRNSummariesBFour{} refreshes, data-change still leads by \RNumAddThreeXOneRNBFourNewTopic{}; the other \RNumAddThreeXThreeDCBFourNewTopicAbsMean{} comes from the one refresh that probing (\RNumRNOverheadShareBFour{} of $B$) takes away. The probe's estimate correlates less strongly with lift-in neighbor counts than $N_{\text{new}}(s)$ on average in every centroid and eight-segment condition.\looseness=-1

For coordinator-routed vector stores in our tested regime ($S = 32, m = 2, c = 16$, $\approx {}$\RNumSteadyEpochInsertVecs{} inserts/epoch), these results suggest three rules accounting for uncharged costs. (1) Under centroid placement at $Q \le 200$, spend bytes on causal online extra search if shards can absorb up to \RNumExtraSearchesBFour{} more HNSW searches/epoch (\RNumESOQFiftyMaxSearches{} at $Q = 50$) and \RNumThreeShards{}-shard fan-out (or an extra round trip). Alternatively, spend them on vector push if the coordinator can buffer and scan recent vectors (\RNumPVMaxVecsBFour{} vectors at $B = 4$). Reserve summary refresh for higher query rates under topic drift, where data-change refresh overtook extra search at $Q = 800$ (fragilely once its score bytes are charged) but trailed vector push. (2) Under arrival-time placement without schedule-aware routing, refresh the active shard summary when arrivals cluster by topic into one segment. On new-topic with $G = 8$ active segments, fresh summaries and vector push are close at $B \le 4$ and vector push leads at $B = 8$. Avoid rebuilding \RNumCentersPerSummary{}-center $k$-means summaries on shards mixing all topics uniformly; there, vector push recovers recent neighbors while leaving the setup summaries, and so routing, unchanged. (3) When scheduling refreshes across active shards, rank shards by the data-change score rather than paying \RNumRNOverheadEpochBytes{} bytes/epoch (plus setup broadcast) for \RNumSampleQueries{}-query cutoff probing. With its \RNumDCScoreEpochBytes{} bytes/epoch charged, data-change still beat the probe at equal refreshes on new-topic at $B \in \{4, 8\}$. The new-vector count $N_{\text{new}}(s)$ tracked lift-in neighbor counts better than the probe in every centroid and eight-segment condition, though it was not tested as a refresh rule. Because summary size is fixed while PV's buffered arrival share at fixed $B$ falls as inserts grow, higher ingest rates may favor refresh (not tested). Beyond our single-process simulator on \RNumCorpusTotalPassages{} MS MARCO MiniLM vectors, future work should test adaptive summaries, schedule-aware routing, higher ingest rates, and tail latency.\looseness=-1
